\section{Why Yaound\'e Communication Is Linguistically Complex}\label{sec:complexity}

The marking scheme asks for a discussion of this question. What follows is argued from the
measurements in this report rather than from general impressions, because the analyzer
happens to quantify exactly the properties that make the problem hard.

\subsection{The lexicon is not drawn from one language}

The \LexiconTotal\ entries carry \OriginCount\ distinct donor-language labels: French,
English, Pidgin, Duala, Beti, Bassa, Bamileke, Fulfulde, Hausa, Bulu, Lingala, German,
Spanish, mixed coinages, and entries whose origin remains genuinely uncertain
(Section~\ref{sec:lexical}). A conventional compiler front end assumes one source language
with one keyword table. Here the keyword table is a \emph{union} of vocabularies that share
no common orthography, and no single dictionary of any one language would recognise even a
majority of a typical utterance.

\subsection{Code-mixing happens inside the sentence, not between sentences}

Speakers do not finish a French sentence and begin an English one. They switch
mid-clause --- a French determiner attached to an English noun governed by a Pidgin verb.
This defeats the obvious engineering shortcut of detecting the language first and then
running a monolingual analyzer, because there is no span of text over which a single
language holds. Our response was to make the donor label a property of the \emph{token}
rather than of the statement, which is why the analyzer can report a mixing ratio at all.

\subsection{Orthography is unstable because the language is primarily spoken}

Francanglais has no standard written form. The same word appears with and without accents,
with French or English spelling conventions, and with variable casing --- and because it is
typed on phones, a typographic apostrophe and an ASCII apostrophe are used
interchangeably, as are precomposed and decomposed Unicode accents that look identical on
screen. The gap between the raw-form and canonical-form frequency tables in
Section~\ref{sec:frequency} is a direct measurement of this instability.

The consequence for the analyzer is that exact string matching is insufficient, which is
why the lexer needs the folded fallback tier. But folding cannot simply be applied
everywhere, because in French it destroys real meaning distinctions:

\begin{center}
\begin{tabular}{@{}llll@{}}
\toprule
\textbf{Pair} & \textbf{Reading A} & \textbf{Reading B} & \textbf{Collapses under folding?} \\
\midrule
\fw{a} / \fw{à}     & \term{VERB} \textit{has}  & \term{PREP} \textit{to}     & yes \\
\fw{la} / \fw{là}   & \term{DET} \textit{the}   & \term{ADV} \textit{there}   & yes \\
\fw{ou} / \fw{où}   & \term{CONJ} \textit{or}   & \term{ADV} \textit{where}   & yes \\
\fw{mais} / \fw{maïs} & \term{CONJ} \textit{but} & \term{NOUN} \textit{maize} & yes \\
\bottomrule
\end{tabular}
\end{center}

All four pairs differ \emph{only} by a diacritic, and each pair spans two different
terminal classes --- so collapsing them changes the token stream the parser sees, not merely
the gloss. The design has to be tolerant of accent loss and simultaneously faithful to
accent meaning. Resolving this tension is the reason for the strict lookup ordering in
Section~\ref{sec:lexical}: exact matches always win, folding is consulted only on failure,
and when folding resolves an ambiguous form the user is told both readings exist.

\subsection{Syntax is looser than written French}

Four constructions in the corpus are ungrammatical in standard written French but entirely
ordinary in speech:

\begin{itemize}[leftmargin=1.4em, itemsep=0.2em]
  \item \textbf{Clause juxtaposition without a conjunction.} Two independent clauses are
        joined by a comma, or by nothing at all.
  \item \textbf{Serial verbs.} Verbs chain directly with no coordination or complementiser.
  \item \textbf{Subject dropping.} A bare predicate stands as a complete utterance, so the
        grammar cannot require an \nonterm{NP} before the verb.
  \item \textbf{Left dislocation.} A topic noun phrase is fronted and then resumed by a
        pronoun, producing two noun phrases before the verb.
\end{itemize}

Each of these forced a specific production into the grammar --- \nonterm{Seq} with \term{SEP}
alternatives, \nonterm{VerbTail}, \nonterm{Clause} $\rightarrow$ \nonterm{Predicate}, and
\nonterm{ClauseTail} $\rightarrow$ \nonterm{NP} \nonterm{Predicate}. A grammar written from
a French textbook would reject all four.

\subsection{Frozen multiword expressions}

\LexiconMultiword\ expressions are lexically frozen: \fw{est-ce que}, \fw{on va faire
comment}. Their meaning is not composed from their parts, so parsing them
compositionally is both wasteful and wrong. They must be recognised at the lexical level as
single \term{FORMULA} tokens --- which means the lexer needs longest-match lookahead over
word sequences, a capability a simple word-at-a-time scanner does not have. It also means
the lexer must know where clause boundaries are, so that a comma cannot be swallowed into
the middle of a formula.

\subsection{Discourse particles carry structure}

Interjections such as \fw{ehh}, \fw{ashia} and \fw{wah} open and close a large share of real
utterances and can constitute a complete turn on their own. In written-language grammars
interjections are usually discarded as noise. Here they are a terminal (\term{INTJ}) with
their own productions, because discarding them would misparse the utterances they frame.

\subsection{Summary}

\begin{keypoint}[The core difficulty]
Yaound\'e urban speech is hard to analyse because none of the simplifying assumptions a
compiler front end normally relies on hold: there is not one source language, there is not
one orthography, there is no authoritative dictionary, and the syntax is more permissive
than any written standard. The engineering answer is not to force the data into a cleaner
model, but to build a pipeline that records provenance at every level --- donor language per
token, evidence per lexicon entry, transformation per grammar rule, and production
identifier per parse step --- so that every claim the analyzer makes can be traced back to
the observation that justifies it.
\end{keypoint}
